\section{A two-wall classification on a nonsymmetric slice}
\label{sec:two-wall-v167}
The determinant construction can be evaluated without listing its
maximal minors. We do so on a slice on which both a jet direction
and a conic direction vary. Put
\[
 S=\Spec\C[b,c]\longrightarrow B_2,\qquad
                  (f,g,r)=(x^2,b,cx).
\]
Let $T$ be the schematic closure of $S_{b\ne0}$ in
$\Gamma_2\times_{B_2}S$. Thus $T$ is the horizontal main component
of this base change, not its full pullback. All target coordinates
$[F:G:R]$ in this section are fixed.

\begin{theorem}[The two-wall surface]\label{thm:two-wall-v167}
The surface $T$ is the blow-up of $S$ at $(b,c)=(0,0)$ followed
by the blow-up of the point $e=c=0$ in the chart $b=ec$.
Equivalently,
\[
                   T\simeq\operatorname{Bl}_{(b,c)(b,c^2)}S.
\]
The maximal-minor ideal in Theorem~\ref{thm:determinantal-v167}
restricts to the explicitly factored ideal
\begin{equation}\label{eq:slice-minors-v167}
          \mathfrak a_2\OO_S=b^{10}(b,c)^4(b,c^2)^6.
\end{equation}
The surface is smooth and has the following three affine charts:
\begin{equation}\label{eq:three-charts-v167}
 \begin{array}{c|c|c}
 \text{chart}&\text{coordinates}&(b,c)\\ \hline
 T_0&(b,k)&(b,bk)\\
 T_1&(e,h)&(e^2h,eh)\\
 T_2&(c,z)&(c^2z,c).
 \end{array}
\end{equation}
On $T_1$ and $T_2$ the universal curve has respectively the
bihomogeneous ideals
\begin{align}
 I_1&=(R^2-hFG,\ tG-esR,\ tR-ehsF),\label{eq:slice-I1-v167}\\
 I_2&=(FG-zR^2,\ tG-czsR,\ tR-csF).\label{eq:slice-I2-v167}
\end{align}
On $T_0$ its ideal sheaf is generated by
\begin{equation}\label{eq:slice-I0-v167}
 t^2G-bs^2F,\quad sR-ktG,\quad tR-bksF,\quad
                         R^2-bk^2FG.
\end{equation}
These equations define flat families of polynomial $3l+1$.
The fibre $T_{(0,0)}$ is a reduced chain $E_{11}\cup E_{21}$
of two projective lines. Their self-intersections are $-2$ and
$-1$, respectively. The first line parametrizes primitive tails,
and the second the conic pencil through the attaching point;
their intersection parametrizes the doubled conic $R^2=0$.
\end{theorem}

\begin{proof}
Perform the stated blow-ups. The first chart of the first blow-up
is $c=bk$ and is unchanged by the second blow-up. In the other
chart write $b=ec$ and blow up $(e,c)$. Its charts $c=eh$ and
$e=cz$ give exactly \eqref{eq:three-charts-v167}. The transition
maps are
\begin{align*}
 k=e^{-1},\quad b=e^2h,
 &\qquad e\ne0 \text{ on }T_1,\ k\ne0\text{ on }T_0,\\
 z=h^{-1},\quad c=eh,
 &\qquad h\ne0 \text{ on }T_1,\ z\ne0\text{ on }T_2.
\end{align*}
In the first overlap the inverse is $e=k^{-1}$, $h=bk^2$;
in the second it is $h=z^{-1}$, $e=cz$. The third overlap is
the composite, so the cocycle identity holds as an identity of
rings. To check the alternative blow-up description, first blow
up $(b,c)$; on the chart $b=ec$ the transform of $(b,c^2)$ is
$c(e,c)$, whereas on $c=bk$ it is the invertible ideal $(b)$.
Removing the invertible factors and applying the universal
property of a blow-up proves the product-ideal assertion.

We next construct the morphism to the Hilbert scheme. On $T_1$
the Hilbert--Burch chart of
Section~\ref{sec:complete-contact-v163} has centre and parameters
$\lambda=B=C=0$, $A=-h$, and retained coefficient $e$.
Its equations are \eqref{eq:slice-I1-v167}. On $T_2$ a
Hilbert--Burch matrix is
\[
 \begin{pmatrix} F&R\\-zR&-G\\-t&-cs\end{pmatrix}.
\]
Its minors are $-(FG-zR^2)$, $tR-csF$, and
$-(tG-czsR)$. On every geometric fibre the ideal has height two:
if $c\ne0$, the equation $tR=csF$ and the other equations give a
one-dimensional graph or its single-contact limit; if $c=0$,
they give the main section and the conic $FG-zR^2$ at $t=0$.
This conic is nonzero for every $z$. There is no common divisor
of the three minors. The fibrewise Hilbert--Burch resolution
\[
 0\longrightarrow\OO(-1,-2)^2\longrightarrow
 \OO(-1,-1)^2\oplus\OO(0,-2)\longrightarrow\OO
 \longrightarrow\OO_C\longrightarrow0
\]
is therefore exact, and its constant Euler characteristic gives
flatness and the polynomial $3l+1$. Equivalently one may apply
the fibrewise flatness criterion successively to this complex of
relatively locally free sheaves.

For $T_0$, on $s\ne0$ eliminate $R$ from
\eqref{eq:slice-I0-v167}. The remaining equation in the relative
$\Pj^1$ with coordinates $[F:G]$ is
\[
                     x^2G=bF,\qquad R=kxG.
\]
It is a relative effective Cartier divisor, since $x^2G-bF$ is
nonzero on every fibre. On $t\ne0$, with $y=s/t$, the equations
instead give $G=by^2F$ and $R=bkyF$, the graph of a section.
These two descriptions prove flatness and exclude any excess
component at infinity. This is also the primitive division family
of Theorem~\ref{thm:division-chart-v162}, with $w=kx$.
On the overlaps substitution of the displayed transitions makes
the ideal sheaves identical; on $T_1\cap T_2$ the quadratic
equations differ by the invertible factor $-h$. Thus the three
flat families glue in the fixed target. Their generic graph is
$[t^2:bs^2:cst]$.

It remains to prove that the constructed surface is the whole
horizontal graph, rather than a resolution mapping to it.
For $a=2$, the matrix $M_2$ has $13$ rows and $75$ columns.
After restricting to $F_0=t^2$, $F_1=bs^2$, $F_2=cst$,
every column has exactly one nonzero entry. Its row index is
$i+2j_0+j_2$ and its entry is $b^{j_1}c^{j_2}$. A nonzero
maximal minor chooses exactly one column from each row.
Consequently its ideal is the product of the ideals generated
by the entries in the individual rows. With $J_1=(b,c)$ and
$J_2=(b,c^2)$ those row ideals, in order from $0$ to $12$, are
\[
 b^4,\ b^3J_1,\ b^2J_2,\ bJ_1J_2,\ J_2^2,\ J_1J_2,
                  \ J_2,\ J_1,\ \OO_S,\OO_S,\OO_S,\OO_S,\OO_S.
\]
They follow by taking all $0\leq i\leq4$ and
$j_0+j_1+j_2=4$ with the indicated row index; divisibility removes
redundant monomials. Their product proves
\eqref{eq:slice-minors-v167}.

The graph closure of the restricted rational Pl\"ucker map is
therefore $\operatorname{Bl}_{b^{10}J_1^4J_2^6}S$, by the graph
argument in Theorem~\ref{thm:determinantal-v167}. This is not an
appeal to unrestricted base change for blow-ups. A principal
nonzero factor does not change the Rees blow-up. Moreover blowing
up $J_1^4J_2^6$ is the same as blowing up $J_1J_2$: on either
integral blow-up, invertibility of the product makes each factor
invertible as a fractional ideal, and conversely. For completeness,
if two nonzero finitely generated ideals have invertible product,
then locally, after dividing their product by its generator, they
are inverse fractional ideals. Thus each is invertible; apply
this repeatedly to the positive powers. The universal properties
supply inverse morphisms between the two blow-ups, identical on
the dense open. This proves that the twice blown-up surface is
exactly $T$. The flat curve constructed above equals its Hilbert
universal curve since the two Hilbert morphisms agree on the
dense open of an integral surface and the Hilbert scheme is
separated. In particular no unproved normality of a graph image
is used in the argument.

On $T_1$ the ideal of the origin pulls back to $(eh)$; on $T_0$
it is $(b)$, and on $T_2$ it is $(c)$. Hence the exceptional fibre
is reduced with a transverse intersection. The strict transform
of the first exceptional curve loses one from its self-intersection
when the second point is blown up. This gives $E_{11}^2=-2$,
$E_{21}^2=-1$, and $E_{11}E_{21}=1$. Substitution $h=0$ and $e=0$
in \eqref{eq:slice-I1-v167} gives its two modular descriptions and
their common doubled conic.
\end{proof}

\begin{corollary}[Complete valued-arc limits on the slice]
\label{cor:slice-chambers-v167}
Let $R$ be a DVR containing $\C$, with uniformizer $\tau$ and
residue field $k$. Let $b=\beta\tau^p$, $c=\gamma\tau^q$, where
$p,q>0$ and $\beta,\gamma\in R^*$; write bars for their residues.
The special embedded curve
has the following ideal sheaf. The displayed bihomogeneous
ideals are understood with irrelevant saturation.
\begin{equation}\label{eq:slice-table-v167}
\begin{array}{c|c|c}
 \text{weights}&\text{residue parameter}&\text{special ideal}\\ \hline
 p<q&k=0&(R,t^2G)\\
 p=q&k=\bar\gamma/\bar\beta&(sR-ktG,tR,R^2)\\
 q<p<2q&& (tG,tR,R^2)\\
 p=2q&h=\bar\gamma^2/\bar\beta&(tG,tR,R^2-hFG)\\
 p>2q&& (tG,tR,FG).
\end{array}
\end{equation}
These five cases exhaust all generically nonincident DVR arcs
centred at the origin with $bc\ne0$ in the fraction field,
including units with arbitrary higher coefficients. If $c=0$
identically, the first row applies. Moreover the determinantal
order on this slice is
\begin{equation}\label{eq:slice-jet-order-v167}
 \mu=10p+4\min(p,q)+6\min(p,2q),
\end{equation}
with $\mu=20p$ when $c=0$.
The two walls $p=q$ and $p=2q$ are actual walls for embedded
limits; coefficients on the walls cannot be omitted.
\end{corollary}
\begin{proof}
For $p\leq q$, use $k=c/b$ on $T_0$. For $q<p\leq2q$ use
$e=b/c$ and $h=c^2/b$ on $T_1$. For $p>2q$ use $z=b/c^2$
on $T_2$. In each case these coordinates belong to $R$, and their
reductions give the table. The formula for $\mu$ follows by taking valuations in
\eqref{eq:slice-minors-v167}; an ideal generated by two elements
has valuation equal to the lesser of their valuations. It is
unaffected by higher unit coefficients. For $k\ne0$ the first generator
$t^2G$ of the primitive central fibre follows from
$sR-ktG$ and $tR$ after taking ideal sheaves. For $k=0$ it must
be retained. This explains why setting $k=0$ in just the last
three generators would give the wrong scheme.
\end{proof}

\begin{proposition}[Components, multiplicities, and the parameter boundary]
\label{prop:slice-geometry-v167}
Every curve in \eqref{eq:slice-table-v167} is Cohen--Macaulay and
has no embedded associated points. In the first three rows the
minimal primes are $(G,R)$ and $(t,R)$; the main section has
multiplicity one and the tail multiplicity two. In the fourth
row the primes are $(G,R)$ and $(t,R^2-hFG)$, each of multiplicity
one. In the fifth they are $(G,R)$, $(t,F)$, and $(t,G)$,
each of multiplicity one. The fourth curve is a main section
attached to a smooth conic; the fifth is a chain of three reduced
lines. The normalization of each reduced curve is the disjoint
union of its smooth components. None of these attached curves
is normal as a whole.

The parameter surface $T$ is normal and smooth. Its exceptional
boundary has simple normal crossings, with
\[
 \begin{array}{c|cc} &E_{11}&E_{21}\\ \hline
 \operatorname{ord}(b)&1&2\\
 \operatorname{ord}(c)&1&1
 \end{array},\qquad
 K_{T/S}=E_{11}+2E_{21}.
\]
These are statements about a parameter-space modification, not a
stable reduction theorem for curves.
\end{proposition}
\begin{proof}
For the primitive rows, the affine presentation
$R=kxG$, $x^2G=0$ is a hypersurface in a smooth relative surface.
Near infinity it is the smooth main section. The two generic
lengths are read from $x^2G$; the residue $k$ changes the
embedding but not the lengths. For the conic rows the
Hilbert--Burch resolution above makes every local ring on the
curve Cohen--Macaulay of dimension one. At the attaching point
use $F=s=1$: the ideals are $(xG,xR,R^2-hG)$, with the analogous
$FG$ presentation in the last row. The ideal decompositions
\[
 (tG,tR,Q)=(G,R)\cap(t,Q),\qquad Q\in\{R^2,R^2-hFG,FG\},
\]
as ideal sheaves give the claimed minimal components and lengths.
A one-dimensional Cohen--Macaulay scheme has no embedded points.
When $h\ne0$ the conic is smooth. When $Q=FG$, its two lines
meet at $[0:0:1]$ and only the line $G=0$ contains the attaching
point $[1:0:0]$. This gives the chain description. The assertions
on the normalization of the reductions follow.

The valuations follow directly from $b=e^2h$, $c=eh$:
$E_{11}$ is $h=0$ and $E_{21}$ is $e=0$. For a blow-up of a
smooth surface at a point the relative canonical divisor is its
exceptional curve, as the Jacobian in a chart shows. Pulling the
first exceptional divisor through the second blow-up adds one
more copy of $E_{21}$. This proves the discrepancy formula.
\end{proof}

\begin{corollary}[Pointwise, but not uniform, finite determinacy]
\label{cor:no-uniform-v167}
Even on $B_2$ no integer $n$ has the property that the coefficient
$n$-jet determines every embedded graph limit. In contrast,
Proposition~\ref{prop:jet-v167} gives an explicit finite bound for
each generically nonincident arc.
\end{corollary}
\begin{proof}
Given $n$, choose $M>n$ and two distinct nonzero constants
$\beta_1,\beta_2$. The arcs
\[
 f=x^2,\qquad g=\beta_i\tau^{2M},\qquad r=\tau^M x
 \quad(i=1,2)
\]
agree modulo $\tau^{n+1}$ and are generically base-point free.
Their limits are the distinct fixed-target conics
$R^2-\beta_i^{-1}FG$ attached to the same main section, by the
fourth row of \eqref{eq:slice-table-v167}. This proves the claim.
\end{proof}

\begin{proposition}[No uniform saturation exponent for full pullbacks]
\label{prop:no-saturation-bound-v167}
The uniformizer saturation exponents of the full pulled-back
parameter graph at contact order two are unbounded. This is
compatible with the finite degree-module statement of
Proposition~\ref{prop:smith-v167}.
\end{proposition}
\begin{proof}
In the symmetric family of Theorem~\ref{thm:ramified-v166}, take
$\delta=\tau^m$. Its full vertical torsion is a nonzero ideal
module tensored with $\C[\tau]/(\tau^m)$. On a chart where that
ideal is nonzero it is killed by $\tau^m$ and not by
$\tau^{m-1}$. Thus the least exponent at which uniformizer
saturation stabilizes is at least $m$. Since $m$ is arbitrary,
there is no uniform bound with contact order fixed at two.
This assertion is about the full parameter-family pullback; it
does not replace the embedded-curve jet theorem with a negative
statement.
\end{proof}
